\section{Astronomical telescopes}
The discipline of observational astronomy is fundamentally an exercise in sampling the electromagnetic field to extract physical parameters of distant celestial sources. The telescope is analyzed not merely as an optical magnifier but as a complex energy-collecting aperture and information processor. The conceptualization of a telescope emphasizes its role in gathering information from electromagnetic radiation. This information is encoded in five primary characteristics of photons: their \textit{direction of arrival}, their \textit{intensity} (or \textit{flux}), the \textit{time of their arrival}, their \textit{frequency} (or \textit{energy}) and their \textit{polarization state}. A telescope serves as the first stage in an observational chain designed to decode these parameters.

The spatial information is captured by the telescope's ability to map a direction $(\theta, \phi)$ on the sky to a specific coordinate $(x, y)$ on a detector. The spectral information is revealed by the frequency distribution of the collected photons, while temporal information is found in the variation of flux over time, such as the rapid pulses from a pulsar or the long-term variability of an active galactic nucleus.

\begin{table}[h!]
	\centering
	\begin{tabular}{|>{\raggedright\arraybackslash}p{0.2\textwidth}|>{\raggedright\arraybackslash}p{0.36\textwidth}|>{\raggedright\arraybackslash}p{0.36\textwidth}|}
		\hline
		\textbf{Information Parameter} & \textbf{Physical Manifestation} & \textbf{Instrumental Technique} \\ \hline
		{Direction} & {	Position on celestial sphere} & {Astrometry / Imaging} \\ \hline
		{Intensity} & {Photon count per unit area/time} & {Photometry} \\ \hline
		{Frequency} & {	Energy of individual photons} & {Spectroscopy} \\ \hline
		{Time} & {Variation in arrival rates} & {Time-series Analysis} \\ \hline
		{Polarization} & {Orientation of electric field} & {Polarimetry} \\ \hline
	\end{tabular}
\end{table}
The design of a telescope is optimized based on which of these parameters is the primary target of the scientific investigation. For example, a telescope intended for high-resolution imaging must prioritize the minimization of diffraction and atmospheric distortion, whereas a telescope designed for faint-object photometry must maximize the collecting area, often at the expense of absolute image sharpness.